\documentclass[aps,prl,twocolumn,showpacs,superscriptaddress]{revtex4-2}
\usepackage{amsmath,amssymb,booktabs,xcolor,hyperref,amsthm}
\newtheorem{theorem}{Theorem}
\newtheorem{remark}{Remark}

\begin{document}

\title{Appendix KA4: Formal Proof of the OHC Mode Capacity Constraint
$|H|\leq n$\\
The topological underpinning of the Pauli Exclusion Theorem
and the Noble Gas Theorem}

\author{Michel Robert Cabri\'e}
\affiliation{Independent Researcher, Victoria, Australia}
\email{ZeroFreeParameters@gmail.com}
\date{20 March 2026}

\begin{abstract}
This appendix derives the OHC mode capacity constraint $|H|\leq n$ 
from the Josephson-Bessel structure of the OHC condensate, providing 
the formal foundation for both the Pauli Exclusion Theorem and the 
Noble Gas Theorem of Letter~81. The constraint follows from the 
requirement that the Hopf phase winding $H\times2\pi\alpha_0$ must 
not exceed the $n$-th Josephson phase $\Delta\varphi_n = 2\pi n\alpha_0$.
The crucial distinction between topological impossibility and energetic 
unfavourability is established.
\end{abstract}

\maketitle

\section{Hopf Charge Quantisation}

The Hopf invariant $H$ classifies maps $\eta: S^3\to S^2$ by the 
homotopy group $\pi_3(S^2) = \mathbb{Z}$ (Hopf 1931). This gives 
$H\in\mathbb{Z}$ exactly --- no continuous values are permitted. 
Hopf charge is additive:
\begin{equation}
    H(f\circ g) = H(f) + H(g)
\end{equation}
This is the topological foundation of charge quantisation in BCT 
(Appendix~JH).

\section{OHC Bessel Modes and Josephson Coupling}

The OHC interior condensate $\Psi_{\rm int}$ satisfies the 
Gross-Pitaevskii equation with Josephson boundary condition at 
coupling $\alpha_0 = r_{\rm oct}\times r_{\rm tet}/\pi = 0.0074081$. 
The linearised solutions are spherical Bessel functions 
$j_n(\kappa_0 r/R)$ indexed by $n = 1,2,3,\ldots$

The Josephson phase accumulated across the $n$-th mode:
\begin{equation}
    \Delta\varphi_n = 2\pi n\alpha_0
\end{equation}

\begin{theorem}[OHC Mode Capacity]
The $n$-th OHC Bessel resonance mode supports total Hopf charge 
$|H_{\rm total}|\leq n$.
\end{theorem}

\begin{proof}
A Hopf excitation of charge $H$ in the $n$-th mode requires phase 
winding $H\times2\pi\alpha_0$. The mode can sustain phase winding up 
to $\Delta\varphi_n = 2\pi n\alpha_0$ without exceeding the Josephson 
energy barrier. Therefore $H\times2\pi\alpha_0 \leq 2\pi n\alpha_0$, 
giving $|H|\leq n$. \qed
\end{proof}

\section{Implications}

\begin{center}
\begin{tabular}{lll}
\toprule
Theorem & Uses $|H|\leq n$ via & Conclusion \\
\midrule
Pauli (L81 Thm~1) & $n=1$: $|2|>1$ forbidden & 2 electrons/mode, opp.~spin \\
Noble Gas (L81 Thm~2) & Closed: no mode for $H\neq0$ & Bonding impossible \\
Shell $2n^2$ (L78) & $n^2$ modes $\times2$ orientations & Shell capacity \\
Bond angle (L80) & Lone pair: $r_{\rm tet}(1+r_{\rm tet})=1/8$ & $\cos\theta=-(1-n_{\rm lp}/8)/3$ \\
\bottomrule
\end{tabular}
\end{center}

\section{Topological vs Energetic}

\begin{remark}
The mode capacity $|H|\leq n$ is a hard topological constraint. 
No smooth deformation of the OHC condensate creates a state with 
$|H|>n$ in the $n$-th mode without a topology change requiring 
infinite energy. By contrast, energetic barriers are finite. 
This is why Pauli exclusion admits no exceptions and noble gas 
inertness is absolute: they are topological theorems.
\end{remark}

\section{The Foundation}

All BCT chemistry rests on:
\begin{align}
    &H \in \mathbb{Z} \quad (\pi_3(S^2)=\mathbb{Z})\\
    &|H| \leq n \text{ per mode} \quad\text{(Josephson-Bessel)}
\end{align}
Two integers. One geometry. All of chemistry.

\begin{thebibliography}{99}
\bibitem{AppJH} M.~R.~Cabri\'e, \textit{Appendix JH}, Zenodo (2026),
doi:10.5281/zenodo.18975018
\bibitem{L78} M.~R.~Cabri\'e, \textit{BCT Letter 78}, Zenodo (2026).
\bibitem{L81} M.~R.~Cabri\'e, \textit{BCT Letter 81}, Zenodo (2026).
\end{thebibliography}

\end{document}
